% GENERATED FILE. Edit canonical lessons, then run npm run generate:books.
% canonical-source-hash: fnv1a64:e281526dff55f1f6
% canonical-lessons: ZH-C466-fadai, ZH-C466-lengjing, ZH-W466-jing-still, ZH-C466-xintiao, ZH-C466-weixiao, ZH-C466-pingjing

\chapter{Keeping Calm}
\label{ch:zh-keeping-calm}
\hlchaptermodality{466}

\begin{chapteropening}
\textbf{By the end of this chapter:} \emph{I can say someone is lost in thought, tell someone to calm down, talk about my heartbeat, smiles and peace of mind.}

Complete the last lesson of chapter 466: I can say someone is lost in thought, tell someone to calm down, talk about my heartbeat, smiles and peace of mind.
\end{chapteropening}

\section[fādāi]{fādāi — to stare blankly; to be lost in a daze}

\label{lesson:ZH-C466-fadai}



\begin{quote}
Before the new one: say the Mandarin for pessimistic, then the Mandarin for sure of oneself.
\end{quote}

\subsection*{You'll want to know: fādāi}
\emph{fādāi} — \textquotedblleft{}to stare blankly; to be lost in a daze\textquotedblright{}.

\begin{quote}
\emph{shàngkè de shíhou tā chángcháng fādāi} — In class he's often lost in a daze.
\end{quote}

Its characters come later; for now it is a word for the ear.

\begin{cousinweb}
\emph{fādāi} is built from \zh{发} 'to set going' and \emph{dāi} 'dull'. The first is inherited from Middle Chinese \emph{pjot} (Guangyun, in Baxter's transcription), glossed 'to start'. The dictionary has no entry under the spelling of the second; under a variant spelling it records \emph{ngoj}, glossed 'foolish', which would regularly give \emph{ái}, not \emph{dāi}, so the history of \emph{dāi} is not settled.
\end{cousinweb}

\begin{grammarlens}[title={What you've built}]
One more word for saying how it was.
\end{grammarlens}

\subsection*{Your turn}
\begin{itemize}
\raggedright
  \item \emph{Say it:} \emph{fādāi}
  \item \emph{Say it:} \emph{fādāi}, once more
  \item \emph{Say it:} \emph{shàngkè de shíhou tā chángcháng fādāi} — In class he's often lost in a daze.
  \item \emph{Recall:} say the Mandarin for pessimistic, then the Mandarin for sure of oneself, then say \emph{fādāi} again
\end{itemize}

\begin{tcolorbox}[breakable,colback=teal!4,colframe=teal!35!black,title={Before you move on}]
What does fādāi mean? (\textquotedblleft{}To stare blankly; to be lost in a daze\textquotedblright{}.) Say it once more.
\end{tcolorbox}
\section[\texorpdfstring{\zh{冷静} (lěngjìng)}{冷静 (lěngjìng)}]{\zh{冷静} (lěngjìng) — calm; cool-headed}

\label{lesson:ZH-C466-lengjing}



\begin{quote}
Before the new one: say the Mandarin for sure of oneself, then the Mandarin for to be lost in a daze.
\end{quote}

\subsection*{You'll want to know: \zh{冷静}}
\textbf{\zh{冷静}} — \emph{lěngjìng} — \textquotedblleft{}calm; cool-headed\textquotedblright{}.

\begin{quote}
\textbf{\zh{你太高兴了先冷静一点儿}} \emph{nǐ tài gāoxìng le xiān lěngjìng yìdiǎnr} — You're getting carried away. Calm down a little first.
\end{quote}

The new character, in pieces: \zh{静} — fourteen strokes in two pieces: on the left the right piece of \zh{情}, a top of four strokes over \zh{月} \emph{yuè}; on the right a short stroke falling to the left, a bar that turns into a stroke falling to the left, a stroke that runs right and turns down, two bars and a vertical with a hook. The next lesson puts it in your hand.

\begin{cousinweb}
\textbf{\zh{冷静}} (\emph{lěngjìng}) is built from \zh{冷} 'cold' and \zh{静} 'still', both inherited from Middle Chinese, \emph{lengX} and \emph{dzjengX} (Guangyun, in Baxter's transcription), which glosses the first 'cold' and the second 'at peace'. The dictionary lists \zh{冷} under three readings, two of them glossed 'cold'; \emph{lengX} is one of those two.
\end{cousinweb}

\begin{grammarlens}[title={What you've built}]
One more word for saying how it was.
\end{grammarlens}

\subsection*{Your turn}
\begin{itemize}
\raggedright
  \item \emph{Say it:} \emph{lěngjìng}
  \item \emph{Say it:} \emph{lěngjìng}, once more
  \item \emph{Say it:} \emph{nǐ tài gāoxìng le xiān lěngjìng yìdiǎnr} — You're getting carried away. Calm down a little first.
  \item \emph{Recall:} say the Mandarin for sure of oneself, then the Mandarin for to be lost in a daze, then say \emph{lěngjìng} again
\end{itemize}

\begin{tcolorbox}[breakable,colback=teal!4,colframe=teal!35!black,title={Before you move on}]
What does \zh{冷静} mean? (\textquotedblleft{}Calm; cool-headed\textquotedblright{}.) Say it once more.
\end{tcolorbox}
\section[\texorpdfstring{\zh{静} (jìng)}{静 (jìng)}]{\zh{静} — the right piece of \zh{情}, then a right side of six}

\label{lesson:ZH-W466-jing-still}



\begin{quote}
The word first, then the character inside it.

\begin{itemize}
\raggedright
  \item \emph{Recall:} say \emph{lěngjìng} — \textbf{R1}, one lesson back, before it decays
\end{itemize}

\begin{quote}
\zh{情} \zh{冷}
\end{quote}

\emph{Read it:} each one aloud, with nothing said first — \emph{qíng}, \emph{lěng}
\end{quote}

\subsection*{Script — the right piece of \zh{情}, then a right side of six}
\hlblockfigure{\detokenize{figures/ZH-W466-jing-still-filmstrip.pdf}}{How \zh{静} is written, stroke by stroke}

\begin{quote}
\zh{静}
\end{quote}

\textbf{Two pieces, left before right.} Count the pieces, not the strokes.

\textbf{The left piece, eight strokes: the right piece of \zh{情}.} A bar, a second bar, a vertical down through them and a long bar; then \zh{月} under them: the stroke down the left, the stroke that runs right, turns down and hooks, and the two short bars inside.

\textbf{The right piece, six strokes.} A short stroke falling to the left. A bar that turns into a stroke falling to the left, like a small roof. A stroke that runs right and turns down. A long bar through it, a second bar under it, then a vertical down the middle that hooks to the left at the foot.

8 + 6 = fourteen.

\subsection*{Your turn}
\begin{itemize}
\raggedright
  \item \emph{Write it:} \zh{静}, all fourteen, with the model in front of you
  \item \emph{Cover:} write it again from memory, and count the pieces, not the strokes
\end{itemize}

\begin{tcolorbox}[breakable,colback=teal!4,colframe=teal!35!black,title={Before you move on}]
How many pieces does \zh{静} have? (\textbf{Two: eight strokes and six}.) Where have you written the left piece before? (\textbf{On the right of \zh{情}}.)

Source: \href{https://raw.githubusercontent.com/chanind/hanzi-writer-data/68d10a4b21150cae5e1ebbd223eed289cf32d90c/data/\%E9\%9D\%99.json}{Hanzi Writer Data \zh{静}, pinned stroke paths and medians}.
\end{tcolorbox}
\section[xīntiào]{xīntiào — heartbeat}

\label{lesson:ZH-C466-xintiao}



\begin{quote}
Before the new one: say the Mandarin for to be lost in a daze, then the Mandarin for cool-headed.
\end{quote}

\subsection*{You'll want to know: xīntiào}
\emph{xīntiào} — \textquotedblleft{}heartbeat\textquotedblright{}.

\begin{quote}
\emph{wǒ gǎnjué xīntiào hěn kuài} — I can feel my heart racing.
\end{quote}

Its characters come later; for now it is a word for the ear.

\begin{cousinweb}
\emph{xīntiào} is built from \zh{心} 'heart' and \emph{tiào} 'to jump', both inherited from Middle Chinese, \emph{sim} and \emph{dew} (Guangyun, in Baxter's transcription), which glosses the second 'to leap'. The regular change from \emph{dew} gives a rising tone; how \emph{tiào} came to its falling tone is not settled.
\end{cousinweb}

\begin{grammarlens}[title={What you've built}]
One more word for saying how it was.
\end{grammarlens}

\subsection*{Your turn}
\begin{itemize}
\raggedright
  \item \emph{Say it:} \emph{xīntiào}
  \item \emph{Say it:} \emph{xīntiào}, once more
  \item \emph{Say it:} \emph{wǒ gǎnjué xīntiào hěn kuài} — I can feel my heart racing.
  \item \emph{Recall:} say the Mandarin for to be lost in a daze, then the Mandarin for cool-headed, then say \emph{xīntiào} again
\end{itemize}

\begin{tcolorbox}[breakable,colback=teal!4,colframe=teal!35!black,title={Before you move on}]
What does xīntiào mean? (\textquotedblleft{}Heartbeat\textquotedblright{}.) Say it once more.
\end{tcolorbox}
\section[wēixiào]{wēixiào — to give a gentle smile}

\label{lesson:ZH-C466-weixiao}



\begin{quote}
Before the new one: say the Mandarin for cool-headed, then the Mandarin for a heartbeat.
\end{quote}

\subsection*{You'll want to know: wēixiào}
\emph{wēixiào} — \textquotedblleft{}to give a gentle smile\textquotedblright{}.

\begin{quote}
\emph{tā wēixiào de shíhou zuì hǎokàn} — She looks her best when she smiles.
\end{quote}

Its characters come later; for now it is a word for the ear.

\begin{cousinweb}
\emph{wēixiào} is built from \emph{wēi} 'slight' and \zh{笑} 'to laugh', both inherited from Middle Chinese, \emph{mj+j} and \emph{sjewH} (Guangyun, in Baxter's transcription), which glosses the first 'fine' and 'small' and the second 'glad'.
\end{cousinweb}

\begin{grammarlens}[title={What you've built}]
One more word for saying how it was.
\end{grammarlens}

\subsection*{Your turn}
\begin{itemize}
\raggedright
  \item \emph{Say it:} \emph{wēixiào}
  \item \emph{Say it:} \emph{wēixiào}, once more
  \item \emph{Say it:} \emph{tā wēixiào de shíhou zuì hǎokàn} — She looks her best when she smiles.
  \item \emph{Recall:} say the Mandarin for cool-headed, then the Mandarin for a heartbeat, then say \emph{wēixiào} again
\end{itemize}

\begin{tcolorbox}[breakable,colback=teal!4,colframe=teal!35!black,title={Before you move on}]
What does wēixiào mean? (\textquotedblleft{}To give a gentle smile\textquotedblright{}.) Say it once more.
\end{tcolorbox}
\section[píngjìng]{píngjìng — calm; peaceful inside}

\label{lesson:ZH-C466-pingjing}



\begin{quote}
Before the new one: say the Mandarin for a heartbeat, then the Mandarin for to give a gentle smile.
\end{quote}

\subsection*{You'll want to know: píngjìng}
\emph{píngjìng} — \textquotedblleft{}calm; peaceful inside\textquotedblright{}.

\begin{quote}
\emph{tīng le yīnyuè yǐhòu wǒ de xīnqíng hěn píngjìng} — After listening to some music, I feel at peace.
\end{quote}

Its characters come later; for now it is a word for the ear.

\begin{cousinweb}
\emph{píngjìng} is built from \emph{píng} 'level' and \zh{静} 'still', both inherited from Middle Chinese, \emph{bjaeng} and \emph{dzjengX} (Guangyun, in Baxter's transcription), which glosses the first 'straight' and 'calm' and the second 'at peace'.
\end{cousinweb}

\begin{grammarlens}[title={What you've built}]
One more word for saying how it was.
\end{grammarlens}

\subsection*{Your turn}
\begin{itemize}
\raggedright
  \item \emph{Say it:} \emph{píngjìng}
  \item \emph{Say it:} \emph{píngjìng}, once more
  \item \emph{Say it:} \emph{tīng le yīnyuè yǐhòu wǒ de xīnqíng hěn píngjìng} — After listening to some music, I feel at peace.
  \item \emph{Recall:} say the Mandarin for a heartbeat, then the Mandarin for to give a gentle smile, then say \emph{píngjìng} again
\end{itemize}

\begin{tcolorbox}[breakable,colback=teal!4,colframe=teal!35!black,title={Before you move on}]
What does píngjìng mean? (\textquotedblleft{}Calm; peaceful inside\textquotedblright{}.) Say it once more.
\end{tcolorbox}
